\section{The deterministic life cycle}\label{sec:det}

Between Diamond's two periods and the stochastic life cycle sits the case with many periods and no
earnings risk, the Auerbach--Kotlikoff class. It is worth doing on its own, because there the
sandwich of Theorem~\ref{thm:sandwich} closes and Light's Theorem 1 becomes algebra.

\subsection{The sandwich closes}

Make income the same in every state, and suppose the household is interior at every age: it saves
something and stops short of the artefactual cap.

\begin{theorem}[Consumption in closed form]\label{thm:detclosed}
Under those two conditions, $c_k(a,z)=\kappa_k\bigl(m+H_k\bigr)$ exactly, so saving is exactly
$(1-\kappa_k)(y+Ra)-\kappa_kH_k$.
\formal{IncomeFluctuation.stageConsumption\_eq\_of\_interior}
\end{theorem}

The proof squeezes the Euler equation between its two directions. Interiority gives both the
inequality above the floor and the one under the cap, hence equality; constant income makes
tomorrow's consumption the same in every state, so the expectation collapses to a point and the
power can be inverted; the budget identity closes the induction. What the stochastic model can
only bracket is here an identity.

\subsection{Uniqueness at logarithmic utility}

\begin{theorem}[Theorem 1, deterministic and logarithmic]\label{thm:detlog}
At $\gamma=1$, saving rises with the interest rate at every age and every asset level, with no
condition. Aggregate capital rises with it, and the stationary equilibrium rate is unique.
\formal{IncomeFluctuation.stagePolicy\_mono\_det\_log},
\lean{IncomeFluctuation.olgCapital\_mono\_det\_log}
\end{theorem}

Two facts do it. At $\gamma=1$ the patience factor is $\beta R$, so the $R$ cancels out of the
recursion for $\kappa$ and the propensity does not depend on the rate at all
(\lean{stageMPC\_one\_eq\_of\_rate}); and human wealth falls as the rate rises
(\lean{stageHumanWealth\_antitone\_rate}). The gap is then exactly
\[
  g(R_2)-g(R_1)=(1-\kappa_k)(R_2-R_1)a+\kappa_k\bigl(H_k(R_1)-H_k(R_2)\bigr),
\]
and both terms are nonnegative. Compare Section~\ref{sec:log}, where at logarithmic utility with
risk the Euler condition is \emph{false} at the rates that matter and a different sufficient
condition has to be found. Risk is what makes it hard, not the horizon.

\subsection{Above logarithmic utility, and a correction}

For $\gamma\ne1$ the propensity moves with the rate and the comparative static becomes a duration
comparison. Consumption at zero assets is $y\sum_{i\le K}R^{-i}\big/\sum_{i\le K}q^i$ with
$q=\Thorn/R$, and it falls with the rate exactly when
\[
  T_y\ \ge\ \Bigl(1-\tfrac1\gamma\Bigr)T_c ,
\]
where $T_y$ and $T_c$ are the durations of the income and consumption profiles. Positive assets
only help, since they add $(1-\kappa_k)(R_2-R_1)a\ge0$.

It is worth being careful about what that threshold is. The certainty-equivalent analysis of
Section~\ref{sec:thm1} reports a critical risk aversion of $4.33$, and it is tempting to read that
as the deterministic threshold. It is not. That number is the certainty equivalent of the
\emph{risky} problem read at the state that matters, where expected income falls back towards the
mean and so is front-loaded relative to consumption. With a flat income profile the two durations
are nearly equal, and the criterion has enormous room (\texttt{numerics/detcheck.m}):

\begingroup\small
\begin{center}
\begin{tabular}{@{}lcccc@{}}
\toprule
$r$ & $0\%$ & $2\%$ & $4\%$ & $6\%$\\
\midrule
Income duration $T_y$ & $29.50$ & $23.70$ & $18.70$ & $14.79$\\
Consumption duration $T_c$ & $28.89$ & $23.40$ & $18.68$ & $14.95$\\
$T_y/T_c$ & $1.021$ & $1.013$ & $1.001$ & $0.989$\\
Criterion binds at $\gamma$ of & never & never & never & $94$\\
\bottomrule
\end{tabular}
\end{center}
\endgroup

Below the rate of time preference the income profile is in fact the \emph{longer} of the two, so
$T_y>T_c$ and the criterion holds for every $\gamma$ whatever; above it consumption is the longer,
and the criterion first fails at a risk aversion near $94$. Either way there is no threshold an
economist would meet: the derivative of saving with respect to the rate at $r=4\%$ is positive and
comfortably so from $\gamma=1$ to $\gamma=5$, at zero assets and at five units of them. The threshold of
Section~\ref{sec:thm1} is a feature of the earnings \emph{profile} the risky problem induces, not
of risk aversion by itself. Only the logarithmic case is formalised, because above it the duration
comparison is an inequality between two geometric sums with different ratios, which is true here
but not unconditionally.
